% Generated draft: requires mathematical and source review.
\documentclass[11pt]{article}
\usepackage[margin=25mm]{geometry}
\usepackage{fontspec}
\setmainfont{Latin Modern Roman}
\usepackage{amsmath,amssymb,mathtools,unicode-math}
\setmathfont{Latin Modern Math}
\usepackage{enumitem,needspace,xurl,hyperref}
\setlength{\parindent}{0pt}
\setlength{\parskip}{5pt}
\setlength{\emergencystretch}{3em}
\newcommand{\N}{\mathbb{N}}
\newcommand{\Z}{\mathbb{Z}}
\newcommand{\Q}{\mathbb{Q}}
\newcommand{\R}{\mathbb{R}}
\newcommand{\C}{\mathbb{C}}
\newcommand{\F}{\mathbb{F}}
\newcommand{\A}{\mathbb{A}}
\newcommand{\PP}{\mathbb{P}}
\newcommand{\E}{\mathbb{E}}
\newcommand{\eps}{\varepsilon}
\newcommand{\dd}{\,\mathrm{d}}
\begin{document}
\section*{UnsolvedMath ID: 700694}
\textbf{Source problem code:} catalogue entry
\subsection{Definitions and mathematical statement}
For a group \(G\) and an element \(u\in G\), the centralizer of \(u\) in \(G\) is the subgroup
\[
C_G(u)=\{g\in G:gu=ug\}.
\]
The supplied problem, qualified as a ``well-known problem,'' asks: describe the structure of \(C_G(u)\) for every almost simple group \(G\) of Lie type and every unipotent element \(u\in G\).

The scope of this statement requires clarification. Under the finite-group interpretation, an almost simple group has a nonabelian finite simple subgroup \(S\) with \(S\leq G\leq\operatorname{Aut}(S)\), identifying \(S\) with its inner automorphism group; ``of Lie type'' would then refer to \(S\). The source does not explicitly specify this interpretation, the underlying fields, or whether the unipotent elements are required to lie in \(S\). For a linear algebraic group in a specified faithful linear realization, an element \(u\) is unipotent when \(u-I\) is nilpotent. The source supplies neither such a realization nor a convention extending unipotence to elements of a general almost simple automorphism extension. Thus the centralizer question is retained with its original terminology, and these unresolved domain qualifications must be settled before it becomes a fully specified classification problem.
\subsection{Short English statement}
Describe the centralizers of unipotent elements in almost simple groups of Lie type. The intended class of groups and the meaning of unipotence in possible automorphism extensions need clarification. The supplied registry labels the problem open; this draft does not independently verify that status.
\subsection{Sources}
\begin{itemize}
\item Local UnsolvedMath registry, permanent ID \(700694\), exact source code \texttt{catalogue entry}, title \emph{Kourovka Notebook Problem 11.1}. The supplied record labels the problem open and records upstream revision \texttt{d371ee917ac9647fbb06cd6fc529deec6f3dda82}.
\item The supplied attribution is to R. Zh. Aleev, proposed in 1990. The background identifies \emph{The Kourovka Notebook (Unsolved Problems in Group Theory)}, 11th issue (1990), original problem number 11.1, and the 21st issue edition (2026), \texttt{arXiv:1401.0300v48}. The problem is described as a ``well-known problem.''
\item The supplied collection description names E. I. Khukhro and V. D. Mazurov as editors of the Kourovka Notebook collection covering Issues 1--21 (1965--2026). These bibliographic details and the open-status label are supplied source information, not independently verified findings.
\end{itemize}
\end{document}
